\documentclass[tikz,border=8pt]{standalone}
\usepackage{ctex}
\usepackage{amsmath}
\usetikzlibrary{calc, arrows.meta}

\begin{document}
\begin{tikzpicture}[scale=1.0, thick]

% 参数：取 alpha=45°, v0已归一化使轨迹适中
% y = x*tan(alpha) - g*x^2/(2*v0^2*cos^2(alpha))
% 取 tan(45)=1, cos^2(45)=0.5, g/(v0^2)=0.5（使最大射程=4）
% y = x - 0.5*x^2/(2*0.5) = x - 0.5*x^2
% 顶点：dy/dx = 1 - x = 0 → x=1, y=0.5
% 落点：y=0 → x(1-0.5x)=0 → x=2（最大射程=2，不够大）
% 改：取 g/(v0^2*cos^2(alpha))=0.4，tan(alpha)=1
% y = x - 0.2*x^2
% 顶点：x=2.5, y=2.5-0.2*6.25=2.5-1.25=1.25
% 落点：x=5

% 坐标轴
\draw[->, thin, gray] (-0.3,0) -- (5.8,0) node[right, font=\small, gray] {$x$};
\draw[->, thin, gray] (0,-0.3) -- (0,2.0) node[above, font=\small, gray] {$y$};
\node[font=\footnotesize, gray] at (-0.25,-0.2) {$O$};

% 弹道抛物线 y = x - 0.2*x^2，x 从 0 到 5
\draw[black, thick, domain=0:5, samples=80, variable=\x]
  plot ({\x}, {\x - 0.2*\x*\x});

% 顶点（最高点）：(2.5, 1.25)
\coordinate (Vtop) at (2.5, 1.25);
\filldraw[black] (Vtop) circle (2pt);
\node[font=\small, above] at (Vtop) {顶点（最高点）};

% 顶点到 x 轴的虚线
\draw[gray, dashed, thin] (Vtop) -- (2.5, 0);
\filldraw[gray] (2.5,0) circle (1.5pt);
\node[font=\footnotesize, gray, below] at (2.5,0) {$x_{\max/2}$};

% 落点（最大射程）：(5, 0)
\coordinate (Land) at (5,0);
\filldraw[black] (Land) circle (2pt);
\node[font=\small, below right] at (Land) {落点};

% 初速度向量箭头（从原点，45°方向）
\coordinate (V0tip) at ({0.9*cos(45)}, {0.9*sin(45)});
\draw[-{Stealth[length=7pt,width=4pt]}, red, very thick]
  (0,0) -- (V0tip) node[above, red, font=\small] {$v_0$};

% 标初始角
\draw[red, thin] (0.35,0) arc[start angle=0, end angle=45, radius=0.35];
\node[red, font=\footnotesize] at (0.52, 0.18) {$\alpha$};

% 原点点
\filldraw[black] (0,0) circle (2pt) node[below left, font=\small] {$O$};

% 轨迹方程
\node[font=\footnotesize, align=left] at (4.5, 1.6) {
  $y=x\tan\alpha - \dfrac{gx^2}{2v_0^2\cos^2\!\alpha}$
};

% 最大射程标注
\draw[gray, thin, <->] (0,-0.45) -- (5,-0.45);
\node[font=\footnotesize, gray, below] at (2.5,-0.45) {最大射程 $R$};

\end{tikzpicture}
\end{document}
